\section{The hidden-formulation study}
\label{sec:hidden-formulation}

This section reports the original 66-case study and keeps it as negative
history. Its broad hypothesis is not retrospectively promoted; the archived,
registered, narrower and powered successor is reported in
Section~\ref{sec:necessity}.

\subsection{Hypotheses}

The design is a prospective comparison rather than a success-by-definition
architecture claim. Its primary hypothesis is:

\begin{quote}
On tasks where the initial formulation, representation, decomposition,
measurement model, or relevant search universe is insufficient, the governed
policy improves root task success versus the strongest matched static or
recursive baseline.
\end{quote}

The safety and mechanistic secondary hypotheses test whether the governed policy
avoids unnecessary reframes on evidence-only and execution-only controls,
reopens dependent closure more precisely than either no reset or full reset, and
preserves mechanic obligations and authority/invariant structure with recursion
depth. The design target for the primary superiority comparison is an absolute
root-success improvement of at least five percentage points; the
unnecessary-reframe rate is prospectively guarded by a two-percentage-point
non-inferiority margin. These are design thresholds, not claimed outcomes.

The registered task families are hidden parent-domain, hidden
representation/coordinate system, hidden decomposition/interface, hidden
measurement/operationalization, evidence-only negative controls, and
execution-only negative controls. Required matched baselines include a static
tool workflow, iterative/tree research, a recursive audit/follow-up controller,
a dependency-aware execution plan, and a revisable information-state controller.
The governed policy is separately ablated by removing explicit $W$, explicit
$M$, typed reframing, dependency-directed reopening, and mechanic self-audit.

The evaluation packet tracks root success, hidden-shift success,
unnecessary-reframe rate, responsibility macro-F1, reframe-target accuracy,
reopen precision/recall/F1, stale-closure survival, invariant/authority
violations, trace fidelity and resource cost. Candidate-caused failures,
timeouts, malformed actions and abstentions remain in the denominator.
Standalone rates receive Wilson intervals; matched system differences use a
paired percentile bootstrap over frozen tasks, with stochastic repetitions
nested within task and system.

\subsection{Prospective design freeze}

The complete machine-readable design freezes hypotheses, task families,
baselines, ablations, metrics, resource and exclusion rules, statistics, access
policy and result-figure/table definitions before final outcomes exist, and
records that no outcome had been accessed at the freeze.

This is deliberately weaker than an execution freeze. The final subject
revision, hidden-shift case-set hash, model and provider versions, baseline
configuration hashes, protected evaluator/adjudication hash, split hashes and
evaluation epoch must all be frozen before the first final test outcome is
inspected; otherwise the run cannot support a publication claim. Any
post-outcome scientific-design change creates a new protocol version rather than
rewriting the frozen one.

\subsection{The local falsifier}

A deterministic suite covers hidden-domain, hidden-representation,
missing-evidence, and execution-only worlds. At the frozen evidence snapshot the
suite passed at a recorded subject commit under a recorded continuous-integration
run.

The important result is not the pass label alone. The suite exposed an
over-broad rule: a singular evidence or execution responsibility could
previously enter the reframe path. That behaviour would have allowed the system
to rewrite the formulation when ordinary evidence acquisition or execution
repair was the supported response. The gate was changed so those cases can no
longer reframe merely because responsibility is singular.

This is architecture-local falsification evidence. It supports the claim that
the implemented contracts distinguish at least these registered known worlds and
that the falsifier was capable of changing the design. It does not establish
open-world adequacy or external superiority.

\subsection{Historical parent-domain replay}

A bounded replay reconstructs an earlier omission in which computational
linguistics and psychometrics were not explicit members of a frozen discipline
basis. Because the historical workflow also required cross-domain and
alternate-vocabulary search, the omission does not prove that faithful execution
could not have discovered those fields. The original retrieval and recognition
trace is not registered, so the historical cause remains only partially
identified.

The current discriminator instead asks whether name-independent functional
signatures can recover mature parent disciplines for scientific-language
interpretation, construct/benchmark validity, and fresh transfer cases. Success
shows that these known worlds can be expressed without inventing a new method
primitive. It does not establish open-world parent-discipline recall.

\subsection{The promotion gate}

The broad, model-bearing publication test was never executed, so its outcome
remains undetermined. Repository tests and protocol prose cannot convert that
state into a superiority claim. The experiment in
Section~\ref{sec:necessity} is an additive successor with new powered worlds,
assimilated parents, execution identities, and a narrower credential-free
mechanical claim; it does not fill or rewrite this gate.

\subsection{Local results}

\subsubsection{Mechanical arm attribution}
The reduction layer, which uses no language model, attributes the correct
top-ranked responsibility on 62 of 66 cases, with 0 of 22 negative controls
receiving a formulation responsibility. Four relations cover 26 cases: a
required input is absent but obtainable (10 cases); a framing remedy was tried
and failed (8 cases); a counterfactual restores the baseline (5 cases); and
equal-weight aggregation is applied over skewed units (3 cases). Half the
attributions rest on rules that fire on exactly one case.

\subsubsection{Solvability audit}
A procedural gold-blind audit rates 55 of 66 cases as solvable from public
content alone. Pre-gold family prediction accuracy is 61 of 66. Five arithmetic
defects were found and confirmed. This is an author-lane case audit, not an
external adjudication or independent naturalistic replication.

\subsubsection{Falsifier repair}
The deterministic suite exposed the over-broad reframe gate described above: a
singular evidence or execution responsibility could previously enter the reframe
path. The gate was repaired so those cases can no longer reframe merely because
responsibility is singular.

\subsubsection{Power and tier status}
The completed archive does not support the broad primary hypothesis. The
governed subject and the strongest registered recursive-audit baseline each
solve 1 of 48 root tasks, for an absolute difference of zero.

The achieved half-width is \AchievedHalfWidth{}, which does not reach the
lowest precision tier the frozen statistical plan registered in advance. The
number of cases required for the primary hypothesis is \RequiredNForHOne{}, and
for the restraint hypothesis \RequiredNForHTwo{}. The study is therefore
underpowered for its primary hypothesis, with \StochasticRepeats{} stochastic
repeats per case.

\subsection{Discussion}

\subsubsection{What the local evidence establishes}
The precursor served as a development falsifier: it exposed the over-broad
reframe gate and forced a repair. Its mechanical reducer maps the top-ranked
responsibility correctly on 62 of 66 authored cases with no formulation false
positives on 22 controls, but that count is not evidence of general
responsibility inference. Thirty-two of the 66 attributions depend on rules
that fire on one case, and a blind responder using only path stems and two fixed
phrases solves 33 of 66 cases. Perfect phrase and filename separation therefore
makes defect discovery---not predictive performance---the scientific output of
this predecessor.

\subsubsection{What the local evidence cannot establish}
The historical $N=\TestCaseCount{}$ falls short of the required
\RequiredNForHOne{} for the primary hypothesis and \RequiredNForHTwo{} for the
restraint hypothesis. The achieved half-width \AchievedHalfWidth{} is below the
lowest registered precision tier. The complete comparison is also adverse:
subject and strongest baseline each solve 1 of 48 root tasks, with difference
zero. The broad H1 is therefore not supported, rather than merely awaiting more
repository checks. This local result authorizes neither external superiority
nor a claim about a wider population. The narrower successor is reported
separately and is never pooled with or used to retroactively validate these
records.

% Table P1-T2: Baseline and ablation comparison (condensed)
% Generated from results/P1-T2_baseline_ablation_results.md
% 2880 records, 12 systems, 48 cases, 5 stochastic repeats

\begin{table}[t]
  \centering
  \caption{Baseline and ablation comparison: root success and unnecessary-reframe rates with 95\%~Wilson score intervals}
  \label{tab:P1-T2}
  \scriptsize
  \setlength{\tabcolsep}{3pt}
  \begin{tabular}{>{\raggedright\arraybackslash}p{3.05cm}lccc}
    \toprule
    System & Role & Root success, & Root success, & Unnecessary \\
    & & all cases & hidden shift & reframe, controls \\
    \midrule
    Governed policy & Subject & 0.0208 [0.0037, 0.1090] & 0.0312 [0.0055, 0.1574] & 0.0000 [0.0000, 0.1936] \\
    Live provider & Baseline & 0.0000 [0.0000, 0.0741] & 0.0000 [0.0000, 0.1072] & 0.0000 [0.0000, 0.1936] \\
    \addlinespace
    \multicolumn{5}{l}{\textit{Baselines}} \\
    Static tool workflow & Baseline & 0.0208 [0.0037, 0.1090] & 0.0312 [0.0055, 0.1574] & 0.0000 [0.0000, 0.1936] \\
    Tree-search research & Baseline & 0.0208 [0.0037, 0.1090] & 0.0312 [0.0055, 0.1574] & 0.1875 [0.0659, 0.4301] \\
    Recursive audit/follow-up & Baseline & 0.0208 [0.0037, 0.1090] & 0.0312 [0.0055, 0.1574] & 0.0000 [0.0000, 0.1936] \\
    Dependency-aware plan & Baseline & 0.0208 [0.0037, 0.1090] & 0.0312 [0.0055, 0.1574] & 0.0000 [0.0000, 0.1936] \\
    Revisable information state & Baseline & 0.0208 [0.0037, 0.1090] & 0.0312 [0.0055, 0.1574] & 0.0000 [0.0000, 0.1936] \\
    \addlinespace
    \multicolumn{5}{l}{\textit{Ablations}} \\
    No explicit $W$ & Ablation & 0.0208 [0.0037, 0.1090] & 0.0312 [0.0055, 0.1574] & 0.0000 [0.0000, 0.1936] \\
    No explicit $M$ & Ablation & 0.0208 [0.0037, 0.1090] & 0.0312 [0.0055, 0.1574] & 0.0000 [0.0000, 0.1936] \\
    Generic retry & Ablation & 0.0208 [0.0037, 0.1090] & 0.0312 [0.0055, 0.1574] & 0.1875 [0.0659, 0.4301] \\
    Full reset & Ablation & 0.0208 [0.0037, 0.1090] & 0.0312 [0.0055, 0.1574] & 0.0000 [0.0000, 0.1936] \\
    No mechanic self-audit & Ablation & 0.0208 [0.0037, 0.1090] & 0.0312 [0.0055, 0.1574] & 0.1875 [0.0659, 0.4301] \\
    \bottomrule
  \end{tabular}
  \par\smallskip
  \textit{Note:} $N=48$ cases per system, 5 stochastic repeats, 2,880 records total; no cell is left undetermined after the live-provider recovery. The recursive audit/follow-up arm is the highest-root-success baseline. All systems except the live-provider arm share one perception layer and are resource-matched. Paired percentile bootstrap (10k resamples) was used for comparisons; Holm correction was applied to secondary families. Full system identifiers, per-family results and mechanistic metrics are archived with the record named in Section~\ref{sec:availability}.
\end{table}


\subsubsection{Failure modes}
The prospective failure taxonomy (Table~\ref{tab:P1-T3}) is populated from the
frozen 2,880-record archive. Abstention dominates across systems (81.4\% of
scored trials), with wrong-responsibility attributions at 9.2\% and missed
reframes at 7.5\%; unnecessary reframes occur only under generic retry, tree
search, and no-self-audit variants. The local suite additionally identified
blind-responder shortcuts in template phrasing, filename-based family
separation, and arithmetic defects in case resources.

% Table P1-T3: Failure taxonomy (roll-up by failure mode, all systems)
% Projected from the archived failure-taxonomy record named in Data and code availability.
\begin{table}[t]
  \centering
  \caption{Failure taxonomy: roll-up by failure mode across all 12 systems and 2,880 scored trials}
  \label{tab:P1-T3}
  \small
  \begin{tabular}{lrrr}
    \toprule
    Failure mode & Trials & Systems & Share \\
    \midrule
    Abstention & 2345 & 11 & 0.8142 \\
    Wrong responsibility & 265 & 4 & 0.0920 \\
    Missed reframe & 215 & 12 & 0.0747 \\
    Unnecessary reframe & 45 & 3 & 0.0156 \\
    Incomplete reopen & 10 & 1 & 0.0035 \\
    \bottomrule
  \end{tabular}
  \par\smallskip
  \textit{Note:} One primary failure mode per trial by precedence (authority~$>$~invariant~$>$~trace integrity~$>$~abstention~$>$~control selectivity~$>$~diagnosis~$>$~targeting~$>$~reopening~$>$~terminal outcome). A lower-priority mechanism can therefore be masked by a higher-priority classification. Abstention dominates across all systems (81.4\%). Wrong-responsibility attributions (9.2\%) and missed reframes (7.5\%) are the next most common; unnecessary reframes occur only under generic retry, tree search, and no-self-audit variants. Per-case detail is in the archived record named in Section~\ref{sec:availability}; the raw suite is never published.
\end{table}


\subsubsection{Limitations of the local suite}
The local falsifier uses constructed cases with frozen semantic properties. It
cannot establish open-world adequacy or external superiority. Template-level
leaks---perfect separation on fixed phrases such as ``The team's framing
is''---were identified and documented but not repaired in this iteration. Half
the mechanical arm attributions rest on singleton rules that may function as
case-specific recognizers.

\subsubsection{Outcome-blind public-source feasibility}
Under a prospectively frozen metadata protocol, the pinned CC0 Retraction Watch
snapshot yielded 49,878 structurally admitted explicit update relations in
42,924 source-connected families. A bounded official Europe PMC rights reducer
identified 12,038 relations in 11,602 families with exact allowlisted licences
on both endpoints. Retraction-notice, correction/erratum, and
expression-of-concern metadata strata each exceeded 20 independent families in
both deterministic waves. Here ``independent'' is only the registered
source-family identity criterion, not a claim of statistical independence.
This establishes feasibility of a rights-valid
three-stratum public development pool, not scientific-action gold, construct
validity, system effect, or external authority.

\subsubsection{Structured-action semantics inside the rights-valid pool}
A separately frozen, outcome-blind successor reconstructed the same 12,038
exact-rights relations and 11,602 families and queried only provider-native
structured metadata. Among 11,991 unique notice records, exactly one of the
three registered strata was assigned to 11,976 records: 11,251 retraction,
256 correction/erratum, and 469 expression of concern. Fifteen records remain
\textsc{Other-or-cannot-check}, and none was forced across multiple strata.
After a prospectively frozen normalization of provider display syntax,
PubMed and Europe PMC relation-type sets agree for 11,988 of 11,991 notices.
The three residual notices are one Europe-PMC-only \texttt{Preprint in}
incidence and two PubMed-only \texttt{RetractionOf} incidences.

This positive result is a construction and interface-conformance result. It
does not identify the action licensed by a scientific episode. Under the
predeclared owner-algebra gate, zero of twelve requirement groups is
sufficient. Provider relation classes offer only partial analogues for a
postpublication coordinate and licensed operations, while item licences offer
only a partial rights analogue; none supplies the required closed-world target
semantics, exhaustive licensed and forbidden operations, recommendation or
execution authority, non-success terminals, target-algebra rights, signed
owner ratification, or independent semantic review. Cross-provider agreement
therefore remains metadata conformance rather than independent gold.

\subsubsection{Public-standard construct feasibility}
A further outcome-blind successor asked whether authoritative standards could
close that authority gap without case text or outcomes. Eight official
documents from four institutional families were byte-bound: NISO CREC,
COPE's retraction guidance, Crossref's Crossmark/policy/relationship
documentation, and NLM JATS relation markup
\cite{nisoCREC2024,copeRetraction2025,crossrefCrossmark2026,nlmJATS14}.
Together they provide a source-native scaffold spanning editorial events,
notices, current-status metadata, update relations, publisher policy, and
relation markup. Under the unchanged conjunctive gate, nine of twelve
owner-algebra groups have a structural analogue, but zero has the required
named-custodian authorship or explicit delegation. The exact nondelegation
upper bound is consequently zero sufficient groups for this library.

This is a positive construct-feasibility result for a rights-bounded,
nonredistributive standards envelope, not a legal determination or a mapping
to an R7 decision. The standards describe or encode postpublication practice;
they do not claim ownership of the R7 vocabulary, its host, or its completed
target corpus. Scientific-action gold therefore remains zero, the
117,649-function audit was not rerun, and its counts remain 0 certified,
116,929 rejected, and 720 \textsc{Cannot-check}.
